\needspace{10\baselineskip}
\begin{center}
\pdfbookmark[1]{2. Zu Art. 53. Projection einer Planetenbahn auf die Ekliptik.}{zusatz-002}
[2. Zu Art. 53.]
\medskip

Projection einer Planetenbahn auf die Ekliptik.
\end{center}
\noindent \(e,\varphi,a,i,\AscendingNode,E\) in gewöhnlicher Bedeutung\\
\(\tilde\omega\) Perihel\\
\(g=\tilde\omega-\AscendingNode\)\\
\(\lambda\) hel[iocentrische] L[än]ge\\
\(\rho\) curtirter Abstand [von der Sonne].

[Für die rechtwinkligen Coordinaten in der Ebene der Ekliptik, gezählt von der Knotenlinie als Abscissenlinie hat man nach Art. 53 und 58 der Theoria motus corporum coelestium]
\[
\begin{aligned}
\text{[1]}\quad &a\cos g(\cos E-e)-a\cos\varphi\sin g\sin E\\
&\qquad=\rho\cos(\lambda-\AscendingNode)\\
&a\cos i\sin g(\cos E-e)+a\cos\varphi\cos i\cos g\sin E\\
&\qquad=\rho\sin(\lambda-\AscendingNode).
\end{aligned}
\]

[Man setze
\[
[2]\quad\begin{aligned}
a\cos g&=\alpha\cos\epsilon\cos\gamma-\beta\sin\epsilon\sin\gamma\\
a\cos\varphi\sin g&=\alpha\sin\epsilon\cos\gamma+\beta\cos\epsilon\sin\gamma\\
a\cos i\sin g&=\alpha\cos\epsilon\sin\gamma+\beta\sin\epsilon\cos\gamma\\
a\cos\varphi\cos i\cos g&=-\alpha\sin\epsilon\sin\gamma+\beta\cos\epsilon\cos\gamma,]
\end{aligned}
\]
\SourceLeaf{298}man bestimme [also] \(\alpha,\beta,\gamma,\epsilon\) durch die Gleichungen
\[
[3]\quad\begin{aligned}
a(1-\cos\varphi\cos i)\cos g&=(\alpha-\beta)\cos(\gamma-\epsilon)\\
a(\cos i-\cos\varphi)\sin g&=(\alpha-\beta)\sin(\gamma-\epsilon)\\
a(1+\cos\varphi\cos i)\cos g&=(\alpha+\beta)\cos(\gamma+\epsilon)\\
a(\cos i+\cos\varphi)\sin g&=(\alpha+\beta)\sin(\gamma+\epsilon),
\end{aligned}
\]
[so wird
\[
\begin{aligned}
&\alpha(\cos(E+\epsilon)-e\cos\epsilon)\cos\gamma\\
&\qquad-\beta(\sin(E+\epsilon)-e\sin\epsilon)\sin\gamma
=\rho\cos(\lambda-\AscendingNode)\\
&\alpha(\cos(E+\epsilon)-e\cos\epsilon)\sin\gamma\\
&\qquad+\beta(\sin(E+\epsilon)-e\sin\epsilon)\cos\gamma
=\rho\sin(\lambda-\AscendingNode)
\end{aligned}
\]
oder]
\[
[4]\quad\begin{aligned}
\alpha\cos(E+\epsilon)-\alpha e\cos\epsilon&=\rho\cos(\lambda-\gamma-\AscendingNode)\\
\beta\sin(E+\epsilon)-\beta e\sin\epsilon&=\rho\sin(\lambda-\gamma-\AscendingNode).
\end{aligned}
\]

Hier sind \(\alpha,\beta\) die beiden halben Hauptaxen der Projectionsellipse;
\(-\alpha e\cos\epsilon,-\beta e\sin\epsilon\) die Coordinaten des Mittelpunkts,
[\(\rho\cos(\lambda-\gamma-\AscendingNode),\allowbreak\rho\sin(\lambda-\gamma-\AscendingNode)\) die rechtwinkligen Coordinaten des Planeten,] die [durch die Sonne zur] Axe \(2\alpha\) [gezogene Parallele] als Abscissenlinie angenommen; \(\gamma+\AscendingNode\) deren Neigung gegen die Äquinoctiallinie.

Bei der Berechnung kann man noch folgende Hülfsgrössen einführen:
\[
\cos i=\operatorname{tang}k,\qquad
\cos\varphi=\operatorname{tang}l,\qquad
\frac{a}{\cos k\cos l}=c;
\]
alsdann ist
\[
[5]\quad\begin{aligned}
a(1-\cos\varphi\cos i)&=c\cos(k+l)\\
a(\cos i-\cos\varphi)&=c\sin(k-l)\\
a(1+\cos\varphi\cos i)&=c\cos(k-l)\\
a(\cos i+\cos\varphi)&=c\sin(k+l).
\end{aligned}
\]

Beispiel. Juno. Elemente für 1810.
\[
\begin{aligned}
\varphi&=14^0\,45'\,35''\\
\log a&=0{,}42648\\
\text{mot. diu. trop.}&=813''{,}5165\quad\text{sid.}=813''{,}379\\
i&=13^0\,4'\,12''\\
\AscendingNode&=171\quad7\quad48\\
\tilde\omega&=53\quad10\quad52
\end{aligned}
\]
\[
\begin{gathered}
\text{Epoche [der mittl. Länge] }1810=95^0\,33'\,50'',\\
g=242^0\,3'\,4''
\end{gathered}
\]
\SourceLeaf{299}
\[
\begin{array}{lcl}
\cos i&.\,.\,.\,.\,.\,.\,.\,.&9{,}988\,6011\\
\cos\varphi&.\,.\,.\,.\,.\,.\,.&9{,}985\,4277\\
\multicolumn{3}{l}{k=44^0\,14'\,53''{,}40}\\
\multicolumn{3}{l}{l=44\quad2\quad20{,}14}\\
\cos k&.\,.\,.\,.\,.\,.\,.&9{,}855\,1097\\
\cos l&.\,.\,.\,.\,.\,.\,.\,.&9{,}856\,6489\\
\hline
&&9{,}711\,7586\\
&&0{,}42648\\
\cline{3-3}
c&.\,.\,.\,.\,.\,.\,.\,.\,.\,.\,.&0{,}71472
\end{array}
\]
\[
\left[\begin{array}{l}k+l=\\k-l=\end{array}\right]
\begin{array}{l}88^0\,17'\,13''{,}54\\0\quad12\quad33{,}26\end{array}
\]
\[
\begin{array}{lcl}
\cos(k+l)&.\,.\,.\,.&8{,}47555\\
\sin(k+l)&.\,.\,.\,.&9{,}99981\\
\cos(k-l)&.\,.\,.\,.&0{,}00000\\
\sin(k-l)&.\,.\,.\,.&7{,}56252\\
\cos g&.\,.\,.\,.\,.\,.\,.\,.&9{,}67088_n\\
\sin g&.\,.\,.\,.\,.\,.\,.\,.&9{,}94614_n
\end{array}
\]
\[
\begin{array}{c}
\left[\begin{array}{l}
(\alpha-\beta)\cos(\gamma-\epsilon)\\
(\alpha-\beta)\sin(\gamma-\epsilon)\\
(\alpha+\beta)\cos(\gamma+\epsilon)\\
(\alpha+\beta)\sin(\gamma+\epsilon)
\end{array}\right]
\begin{array}{cl}
.\,.&8{,}86115_n\\.\,.&8{,}22338_n\\.\,.&0{,}38560_n\\.\,.&0{,}66067_n
\end{array}\\
\hline
\end{array}
\]
\[
\left[\begin{array}{l}\gamma-\epsilon=\\\gamma+\epsilon=\end{array}\right]
\begin{array}{l}192^0\,58'\,2''\\242\quad2\quad26\end{array}
\]
\[
\begin{array}{lcl}
\alpha-\beta&.\,.\,.\,.\,.\,.\,.\,.\,.\,.\,.&8{,}87237\\
\alpha+\beta&.\,.\,.\,.\,.\,.\,.\,.\,.\,.\,.&0{,}71457
\end{array}
\]
\[
\begin{aligned}
\gamma&=217^0\,30'\,14''\\
\epsilon&=24\quad32\quad12\\
\gamma+\AscendingNode&=28\quad38\quad2
\end{aligned}
\]
\[
\begin{array}{lcl}
\alpha&.\,.\,.\,.\,.\,.\,.\,.\,.\,.\,.\,.\,.\,.&0{,}41974\\
\beta&.\,.\,.\,.\,.\,.\,.\,.\,.\,.\,.\,.\,.\,.&0{,}40725\\
e&.\,.\,.\,.\,.\,.\,.\,.\,.\,.\,.\,.\,.\,.&9{,}40614\\
\hline
\end{array}
\]
\[
\begin{array}{c}
\left[\begin{array}{l}
\alpha e\\\cos\epsilon\\\beta e\\\sin\epsilon
\end{array}\right]
\begin{array}{cl}
.\,.\,.\,.\,.\,.\,.\,.\,.\,.&9{,}82588\\
.\,.\,.\,.\,.\,.\,.\,.\,.\,.&9{,}95890\\
.\,.\,.\,.\,.\,.\,.\,.\,.\,.&9{,}81339\\
.\,.\,.\,.\,.\,.\,.\,.\,.\,.&9{,}61834
\end{array}\\
\hline
\end{array}
\]
\[
\left.\begin{aligned}
-\alpha e\cos\epsilon&=-0{,}60923\\
-\beta e\sin\epsilon&=-0{,}27022
\end{aligned}\right\}
\begin{array}{l}\text{Coord.}\\\text{C[entri].}\end{array}
\]

[Man kann also die heliocentrischen Coordinaten jetzt nach der Formel
\[
\begin{aligned}
&\rho\cos(\lambda-28^0\,38'\,2'')\\
&\qquad=-0{,}60923+(0{,}41974)\cos(E+24^0\,32'\,12'')\\
&\rho\sin(\lambda-28\quad38\quad2)\\
&\qquad=-0{,}27022+(0{,}40725)\sin(E+24\quad32\quad12)
\end{aligned}
\]
berechnen, sobald nur die excentrische Anomalie bekannt ist.

Die Zahlen in Parenthese bedeuten Logarithmen.]
\begin{center}\rule{.3\linewidth}{.4pt}\end{center}

Die Hauptmomente der Auflösung der umgekehrten Aufgabe: aus der Projection einer Planetenbahn und ihres Brennpunkts, jene selbst zu finden, sind folgende:

\SourceLeaf{300}Es sei \(\alpha\) die halbe grosse, \(\beta\) die halbe kleine Axe der Projectionsellipse; ferner, vom Mittelpunkt der Ellipse an gerechnet, gegen eine willkürliche Abscissenlinie
\begin{center}
\begin{tabular}{@{}lcl@{}}
\(P\)&.\ .\ .\ .&Richtung von \(\alpha\)\\
\(Q\)&.\ .\ .\ .&\begin{minipage}[t]{.75\linewidth}Richtung der Linie zum Brennpunkte, deren Grösse \(=\delta\).\end{minipage}
\end{tabular}
\end{center}

Man setze
\[
[6]\quad\operatorname{tang}(R-P)\,.\,\operatorname{tang}(Q-P)
=-\frac{\beta\beta}{\alpha\alpha}\mathpunct{\raisebox{.45ex}{,}}
\]
[so dass also \(R\) die Richtung des zur Linie zum Brennpunkt conjugirten Durchmessers ist; wenn \(q\) den durch den Brennpunkt gehenden Halbmesser und \(r\) den ihm conjugirten bezeichnet, so ist]
\[
[7]\quad\begin{aligned}
\BareRoot{\left\{\frac{\cos(Q-P)^2}{\alpha\alpha}+
\frac{\sin(Q-P)^2}{\beta\beta}\right\}}&=\frac{1}{q}\\
\BareRoot{\left\{\frac{\cos(R-P)^2}{\alpha\alpha}+
\frac{\sin(R-P)^2}{\beta\beta}\right\}}&=\frac{1}{r}\mathpunct{\raisebox{.45ex}{,}}
\end{aligned}
\]
also
\[
[8]\quad\frac{qq}{rr}=-\frac{\sin2(R-P)}{\sin2(Q-P)}\mathpunct{\cdot}
\]

[Die Excentricität der Ellipse ist]
\[
[9]\quad e=\frac{\delta}{q}\mathpunct{\cdot}
\]

[Durch Einführung des auch im vorigen benutzten Winkels \(\epsilon\), welcher die zum Radius Vector \(q\) gehörige excentrische Anomalie in der Projectionsellipse darstellt, lassen sich die Endpunkte der conjugirten Radii Vectores \(q\) und \(r\) in rechtwinkligen Coordinaten, jetzt bezogen auf die Axe \(2\alpha\) als Abscissenlinie, wie folgt, ausdrücken:
\[
\begin{aligned}
q\cos(Q-P)&=\alpha\cos\epsilon,&
r\cos(R-P)&=-\alpha\sin\epsilon,\\
q\sin(Q-P)&=\beta\sin\epsilon,&
r\sin(R-P)&=\beta\cos\epsilon;
\end{aligned}
\]
leitet man aus diesen die rechtwinkligen Coordinaten, bezogen auf die Knotenlinie als Abscissenlinie ab, indem man um den Winkel \(\gamma=P-\AscendingNode\) dreht, so hat man mit Rücksicht auf die Gleichungen 2:]
\[
[10]\quad\begin{aligned}
q\cos(Q-\AscendingNode)&=a\cos g,\\
\frac{r}{\BareRoot{(1-ee)}}\cos(R-\AscendingNode)&=-a\sin g\\
q\sin(Q-\AscendingNode)&=a\cos i\sin g,\\
\frac{r}{\BareRoot{(1-ee)}}\sin(R-\AscendingNode)&=a\cos i\cos g,
\end{aligned}
\]
\SourceLeaf{301}[woraus man leicht ableitet]
\[
[11]\quad\frac{qq}{rr}(1-ee)
=-\frac{\sin2(R-\AscendingNode)}{\sin2(Q-\AscendingNode)}\mathpunct{\cdot}
\]

[Ferner findet sich]
\[
[12]\quad\begin{aligned}
qq+\frac{rr}{1-ee}&=aa(1+\cos i^2)\\
\frac{qr}{\BareRoot{(1-ee)}}\sin(R-Q)&=aa\cos i\\
\frac{qr}{\BareRoot{(1-ee)}}\cos(R-Q)&=-\tfrac{1}{2}aa\sin i^2\sin2g\\
qq-\frac{rr}{1-ee}&=aa\sin i^2\cos2g.
\end{aligned}
\]

[In ähnlicher Weise leitet man ab
\[
\begin{aligned}
\frac{qr}{\BareRoot{(1-ee)}}\cos(Q+R-2\AscendingNode)
&=-\tfrac{1}{2}aa(1+\cos i^2)\sin2g\\
\frac{qr}{\BareRoot{(1-ee)}}\sin(Q+R-2\AscendingNode)
&=aa\cos i\cos2g,
\end{aligned}
\]
also mit Rücksicht auf die Gleichungen 12:]
\[
[13]\quad\begin{aligned}
&aa\sin i^2\cos(Q+R-2\AscendingNode)\\
&\qquad=\left(qq+\frac{rr}{1-ee}\right)\cos(R-Q)\\
&aa\sin i^2\sin(Q+R-2\AscendingNode)\\
&\qquad=\left(qq-\frac{rr}{1-ee}\right)\sin(R-Q).
\end{aligned}
\]

[Die abgeleiteten Gleichungen enthalten alles nöthige zur Lösung der Aufgabe und können auf mannigfaltige Weise benutzt werden; sind \(\alpha,\beta,P,Q,\delta\) gegeben, so kann man z.\,B. \(R\) aus 6, \(q\) und \(r\) aus 7, \(e\) aus 9 bestimmen; sodann finden sich \(\AscendingNode,a,i,g\) aus 13 und 10. Die Lage der Abscissenlinie, von der die Richtungen \(P,Q,R,\AscendingNode\) gerechnet werden, bleibt vollkommen willkürlich.]
\begin{center}\rule{.3\linewidth}{.4pt}\end{center}

